\chapter{把排除边界推到二分之一}
\label{chap:finish}
\section{先排除临界线右侧}
现在将引用定理\ref{thm:arithmetic}与已经证明的窗口判据\ref{thm:window-zeta}接起来。

反设 $\rho\ne1$ 是一个零点，且
\[
\beta=\Rea\rho>\frac12.
\]
令
\[
a=\min\{\beta-\tfrac12,\tfrac12\}>0,
\qquad \theta=\frac12+\frac a2.
\]
因此 $0<\theta<1$ 且 $\theta<\beta$。算术定理的常数 $K>0$ 与 $\varepsilon_0>0$ 先固定，再取
\begin{equation}
\varepsilon=\min\left\{\frac{\varepsilon_0}{2},\frac{a}{4K}\right\}.
\label{eq:eps-final}
\end{equation}
这保证 $0<\varepsilon<\varepsilon_0$，并有
\[
\frac12+K\varepsilon\le\frac12+\frac a4
<\frac12+\frac a2=\theta<\beta.
\]
在算术定理提供的任意远窗口上，因 $x\ge1$，
\[
G(x)\le x^{1/2+K\varepsilon}\le x^\theta.
\]
窗口判据于是排除所有实部大于 $\theta$ 的零点，与 $\Rea\rho=\beta>\theta$ 矛盾。我们得到
\begin{equation}
\boxed{\zeta(s)\ne0\quad\text{当 }\Rea s>\tfrac12,\ s\ne1.}
\label{eq:right-free}
\end{equation}

\Expl{不必真的把 $\varepsilon$ 设为零。一个假想右侧零点与临界线之间有固定正距离；只要选择一次足够小但仍为正的 $\varepsilon$，就已经把它排除了。不同假想零点可以对应不同的小参数与窗口。}

\section{把函数方程写成可使用的形式}
定义完备化函数
\[
\mathcal Z(s)=\pi^{-s/2}\Gamma(s/2)\zeta(s).
\]
经典函数方程为亚纯恒等式
\begin{equation}
\mathcal Z(s)=\mathcal Z(1-s).
\label{eq:FE}
\end{equation}
例如，标准的 $\xi(s)=\tfrac12s(s-1)\mathcal Z(s)$ 满足
$\xi(s)=\xi(1-s)$；在 $s\ne0,1$ 处除去相等的多项式因子，再作亚纯延拓，就得到式\eqref{eq:FE}。标准公式见资料[4]。

Gamma 函数在 $0,-1,-2,\ldots$ 有单极点，在其他地方有限且非零。因而 $\Gamma(s/2)$ 的极点恰位于 $s=0,-2,-4,\ldots$。另外，$\pi^{-s/2}=\exp(-s\log\pi/2)$ 始终非零。Gamma 的这些性质见资料[5]。

\section{先核对 \texorpdfstring{$\zeta(0)$}{ζ(0)}，再处理左侧零点}
$\zeta$ 在 $0$ 全纯。由 $\Gamma(s/2)=2/s+O(1)$，
\[
\mathcal Z(s)=\frac{2\zeta(0)}s+O(1)\quad(s\to0).
\]
另一边，$\zeta(1-s)=-1/s+O(1)$，且
$\pi^{-(1-s)/2}\Gamma((1-s)/2)\to\pi^{-1/2}\Gamma(1/2)=1$。
函数方程给出 $2\zeta(0)=-1$，所以
\begin{equation}
\zeta(0)=-\frac12\ne0.
\label{eq:zeta-zero-value}
\end{equation}

现在设 $\rho$ 为非平凡零点且 $\Rea\rho<1/2$。由式\eqref{eq:zeta-zero-value}有 $\rho\ne0$；根据非平凡的限定，$\rho$ 也不是负偶数。因此 $\Gamma(\rho/2)$ 有限非零，$\mathcal Z$ 在 $\rho$ 有意义，并且
\[
\mathcal Z(\rho)=\pi^{-\rho/2}\Gamma(\rho/2)\zeta(\rho)=0.
\]
函数方程于是给出 $\mathcal Z(1-\rho)=0$。
但
\[
\Rea(1-\rho)=1-\Rea\rho>\frac12,
\qquad 1-\rho\ne1.
\]
由右侧非消失式\eqref{eq:right-free}，$\zeta(1-\rho)\ne0$。又因为
$\Rea((1-\rho)/2)>1/4$，Gamma 因子也有限非零，故
\[
\mathcal Z(1-\rho)=
\pi^{-(1-\rho)/2}\Gamma((1-\rho)/2)\zeta(1-\rho)\ne0,
\]
矛盾。左侧非平凡零点也被排除了。

\begin{center}
\begin{tikzpicture}[x=6cm,y=.65cm,>=Latex]
\fill[Teal!9] (.5,-2.5) rectangle (1.05,2.5);
\draw[->] (-.05,0)--(1.13,0) node[right] {$\Rea s$};
\draw[->] (0,-2.8)--(0,2.9) node[above] {$\Ima s$};
\draw[Teal,thick] (.5,-2.5)--(.5,2.5);
\node[below right] at (.5,0) {$1/2$};
\fill[red!65!black] (.25,2) circle (1.5pt);
\fill[red!65!black] (.75,-2) circle (1.5pt);
\draw[->,thick] (.25,2)--(.75,-2);
\node[above] at (.25,2) {$\rho=\beta+it$};
\node[below right] at (.75,-2) {$1-\rho=(1-\beta)-it$};
\node[align=center] at (.82,1.5) {右侧已被排除};
\end{tikzpicture}
\end{center}
图中是用于反证的假想位置。变换 $s\mapsto1-s$ 是绕 $(1/2,0)$ 的半周旋转，虚部也改变符号；这里没有把它画成同高度的横向反射。

\section{得到结论}
非平凡零点的实部既不能大于 $1/2$，也不能小于 $1/2$。所以
\[
\boxed{\Rea\rho=\frac12.}
\]
这完成了定理\ref{thm:RH}在引用定理\ref{thm:arithmetic}及明确列出的标准解析基础上的证明。

